\documentclass[UTF8,11pt]{ctexart}
\usepackage[a4paper,margin=24mm]{geometry}
\usepackage{amsmath,amssymb,amsthm,mathtools,mathrsfs}
\usepackage[all]{xy}
\usepackage{xurl,hyperref,enumitem}
\usepackage{tikz}
\usetikzlibrary{arrows.meta,cd,decorations.markings,calc,patterns}
\hypersetup{hidelinks,breaklinks=true}
\setlength{\emergencystretch}{3em}
\allowdisplaybreaks
\newtheorem*{theorem}{Theorem}
\newtheorem*{thm}{Theorem}
\newtheorem*{lemma}{Lemma}
\newtheorem*{proposition}{Proposition}
\newtheorem*{prop}{Proposition}
\newtheorem*{corollary}{Corollary}
\newtheorem*{cor}{Corollary}
\newtheorem*{claim}{Claim}
\theoremstyle{definition}
\newtheorem*{definition}{Definition}
\newtheorem*{defn}{Definition}
\newtheorem*{remark}{Remark}
\newtheorem*{example}{Example}
\newcommand{\R}{\mathbb R}
\newcommand{\C}{\mathbb C}
\newcommand{\Z}{\mathbb Z}
\newcommand{\Q}{\mathbb Q}
\newcommand{\N}{\mathbb N}
\newcommand{\F}{\mathbb F}
\newcommand{\D}{\mathbb D}
\newcommand{\bB}{\mathbb B}
\newcommand{\bC}{\mathbb C}
\newcommand{\bR}{\mathbb R}
\newcommand{\bZ}{\mathbb Z}
\newcommand{\bN}{\mathbb N}
\newcommand{\bQ}{\mathbb Q}
\newcommand{\bD}{\mathbb D}
\newcommand{\bF}{\mathbb F}
\newcommand{\CP}{\mathbb{CP}}
\newcommand{\RP}{\mathbb{RP}}
\newcommand{\sO}{\mathscr O}
\newcommand{\sI}{\mathscr I}
\newcommand{\sF}{\mathscr F}
\newcommand{\sG}{\mathscr G}
\newcommand{\sH}{\mathscr H}
\newcommand{\sR}{\mathscr R}
\newcommand{\sM}{\mathscr M}
\newcommand{\sV}{\mathscr V}
\newcommand{\sU}{\mathscr U}
\DeclareMathOperator{\Hom}{Hom}
\DeclareMathOperator{\Ext}{Ext}
\DeclareMathOperator{\Tor}{Tor}
\DeclareMathOperator{\Spec}{Spec}
\DeclareMathOperator{\Proj}{Proj}
\DeclareMathOperator{\supp}{supp}
\DeclareMathOperator{\im}{im}
\DeclareMathOperator{\id}{id}
\DeclareMathOperator{\Tr}{Tr}
\DeclareMathOperator{\tr}{tr}
\DeclareMathOperator{\rank}{rank}
\DeclareMathOperator{\ord}{ord}
\DeclareMathOperator{\dist}{dist}
\DeclareMathOperator{\End}{End}
\DeclareMathOperator{\Aut}{Aut}
\DeclareMathOperator{\Gal}{Gal}
\DeclareMathOperator{\vol}{vol}
\DeclareMathOperator{\Cl}{Cl}
\DeclareMathOperator{\coker}{coker}
\DeclareMathOperator{\codim}{codim}
\DeclareMathOperator{\divergence}{div}
\DeclareMathOperator{\Ric}{Ric}
\DeclareMathOperator{\grad}{grad}
\DeclareMathOperator{\Hess}{Hess}
\newcommand{\abs}[1]{\left\lvert#1\right\rvert}
\newcommand{\sabs}[1]{\lvert#1\rvert}
\newcommand{\babs}[1]{\bigl\lvert#1\bigr\rvert}
\newcommand{\Babs}[1]{\Bigl\lvert#1\Bigr\rvert}
\newcommand{\bbabs}[1]{\biggl\lvert#1\biggr\rvert}
\newcommand{\BBabs}[1]{\Biggl\lvert#1\Biggr\rvert}
\newcommand{\norm}[1]{\left\lVert#1\right\rVert}
\newcommand{\snorm}[1]{\lVert#1\rVert}
\newcommand{\bnorm}[1]{\bigl\lVert#1\bigr\rVert}
\newcommand{\Bnorm}[1]{\Bigl\lVert#1\Bigr\rVert}
\newcommand{\bbnorm}[1]{\biggl\lVert#1\biggr\rVert}
\newcommand{\BBnorm}[1]{\Biggl\lVert#1\Biggr\rVert}
\newcommand{\widebar}[1]{\overline{#1}}
\newcommand{\nicefrac}[2]{\frac{#1}{#2}}
\newcommand{\linnprod}[2]{\langle#1,#2\rangle}
\newcommand{\myindex}[1]{#1}
\newcommand{\glsadd}[1]{}
\newcommand{\avoidbreak}{}
\newcommand{\sourceref}[1]{\textnormal{[原文：\nolinkurl{#1}]}}
\newcommand{\thmref}[1]{\sourceref{#1}}
\newcommand{\lemmaref}[1]{\sourceref{#1}}
\newcommand{\propref}[1]{\sourceref{#1}}
\newcommand{\corref}[1]{\sourceref{#1}}
\newcommand{\defnref}[1]{\sourceref{#1}}
\newcommand{\exampleref}[1]{\sourceref{#1}}
\newcommand{\exerciseref}[1]{\sourceref{#1}}
\newcommand{\figureref}[1]{\sourceref{#1}}
\newcommand{\sectionref}[1]{\sourceref{#1}}
\newcommand{\appendixref}[1]{\sourceref{#1}}
\newcommand{\stacksref}[2]{\href{https://stacks.math.columbia.edu/tag/#1}{#2 [#1]}}


\newcommand{\E}{\mathbb E}
\newcommand{\Prob}{\mathbb P}
\newcommand{\Reff}{\mathcal R_{\mathrm{eff}}}
\newcommand{\eps}{\varepsilon}
\DeclareMathOperator{\diam}{diam}
\DeclareMathOperator{\Var}{Var}
\DeclareMathOperator{\Crit}{Crit}
\newcommand{\dd}{\,\mathrm d}

\begin{document}
\section*{053\quad 圆堆积载体与随机游走的常返暂留性}
\subsection*{来源与计数目标}
Asaf Nachmias, \emph{Planar Maps, Random Walks and Circle Packing}, LNM 2243, Springer, 2020，第4章，Theorem 4.4 的第(3)、(4)项，\S4.4.1--4.4.2，印刷页52--53；附 Lemma 4.2 及 Lemma 4.9。实际读取开放获取网页正文并对照全书 PDF 的可读数学文本，2026-09-28。
\par\url{https://link.springer.com/chapter/10.1007/978-3-030-27968-4_4}
\par\textbf{本文件只计一个问题：}给定有界度平面三角剖分的圆堆积，证明载体为平面时常返、载体为单位圆盘时暂留。这是作者 Theorem 4.4 的两个指定分项，不是四项全定理。
\subsection*{作者命题与人类证明}
\paragraph{Notation (Sections 4.2--4.3).}
\noindent\textit{以下背景与记号据所列原文章节摘编；不是逐字引文，也不是新增证明。}\par
For a circle packing $P$ and a vertex $v$, denote by $C_v$ the circle corresponding to $v$, by $\operatorname{cent}(v)$ the center of that circle, and by $\operatorname{rad}(v)$ its radius. We write $G(P)$ for the tangency graph of the packing $P$, that is, the graph in which each vertex is a circle of $P$ and two such circles form an edge when they are tangent.
Given a circle packing $P$ of a triangulation $G$, we define the carrier of $P$, denoted $\operatorname{Carrier}(P)$, to be the union of the closed discs bounded by the circles of $P$ together with the spaces bounded between any three circles that form a face (i.e., the interstices). We say that $G$ is circle packed in $\R^2$ when $\operatorname{Carrier}(P)=\R^2$. Denote by $\mathbb U$ the disk $\{z\in\R^2:|z|<1\}$; we say that $G$ is circle packed in $\mathbb U$ when $\operatorname{Carrier}(P)=\mathbb U$.

\begin{theorem}[Theorem 4.4, parts 3 and 4, He and Schramm]
Let $G$ be an infinite simple plane triangulation with bounded degrees.
\begin{enumerate}[label=(\arabic*),start=3]
\item If $P$ is a circle packing of $G$ with $\operatorname{Carrier}(P)=\R^2$, then $G$ is recurrent.
\item If $P$ is a circle packing of $G$ with $\operatorname{Carrier}(P)=\mathbb U$, then $G$ is transient.
\end{enumerate}
\end{theorem}

\paragraph{Ring Lemma (Lemma 4.2, used below).}
Given circles $C_0,C_1,\ldots,C_M$ with disjoint interiors, we say that $C_1,\ldots,C_M$ completely surround $C_0$ if they are all tangent to $C_0$ and $C_i$ is tangent to $C_{i+1}$ for $i=1,\ldots,M$ (where $C_{M+1}$ is set to be $C_1$).
For every integer $M>0$ there exists $A>0$ such that if $C_0$ is a circle completely surrounded by $M$ circles $C_1,\ldots,C_M$, and $r_i$ is the radius of $C_i$ for every $i=0,1,\ldots,M$, then $r_0/r_i\le A$ for every $i=1,\ldots,M$.
\begin{proof}[Author's proof of Lemma 4.2]
We may scale the picture so that $r_0=1$. Assume that the radius of $C_2$ is small and consider the circles $C_1$ and $C_3$ to its left and right. It cannot be that both $C_1$ and $C_3$ have large radii compared to $C_2$ since in this case they will intersect; see Fig. 4.1. Hence, one of them has to be small as well. Assume without loss of generality that it is $C_3$. By similar reasoning, one of $C_1$ and $C_4$ has to be small. We continue this argument this way and get a path of circles of small radii; thus, for the circles $C_1,\ldots,C_M$ to completely surround $C_0$ we learn that $M$ must be large.
\end{proof}

\paragraph{Section 4.4.1: Proof of part 3.}
Denote the circle packing $P=\{C_v\}_{v\in V}$ where $V$ is the vertex set of $G$ and $C_v$ denotes the circle corresponding to the vertex $v$. Write $\Delta$ for the maximum degree of $G$ and fix a vertex $v_0$. By scaling and translating we may assume that $C_{v_0}$ is a radius $1$ circle around the origin. For a real number $R>0$, let $V_R=V_{B(0,R)}$ denote set of vertices $v$ for which $\operatorname{cent}(v)$ is in the Euclidean ball of radius $R$ around the origin.
\begin{lemma}[4.9]
There exist $C=C(\Delta)>1$ and $c=c(\Delta)>0$ such that for every $R\ge1$ we have
\begin{enumerate}[label=(\roman*)]
\item There are no edges between $V_R$ and $V\setminus V_{CR}$, and
\item $\Reff(V_R\leftrightarrow V\setminus V_{CR})\ge c$.
\end{enumerate}
\end{lemma}
\begin{proof}
We begin with part (i). For every $v\in V_R$ it holds that $\operatorname{rad}(v)\le R$ since $C_{v_0}$ is centered at the origin. By the Ring Lemma (Lemma 4.2), there exists $A=A(\Delta)$ such that $\operatorname{rad}(u)\le AR$ for every $u\sim v$, and therefore $|\operatorname{cent}(u)|\le(A+2)R$. Hence (i) holds with $C=A+2$.
To prove part (ii) we define
\[
h(v)=\begin{cases}
0&v\in V_R,\\1&v\in V\setminus V_{CR},\\
\dfrac{|\operatorname{cent}(v)|-R}{(C-1)R}&\text{otherwise}.
\end{cases}
\]
Recall from Lemma 2.32 that $\Reff(V_R\leftrightarrow V\setminus V_{CR})\ge\mathcal E(h)^{-1}$. By the triangle inequality, for an edge $\{x,y\}$ with both endpoints in $V_{CR}\setminus V_R$ we have
\[
|h(x)-h(y)|\le\frac{|\operatorname{cent}(x)-\operatorname{cent}(y)|}{(C-1)R}
=\frac{\operatorname{rad}(x)+\operatorname{rad}(y)}{(C-1)R},
\]
and it is straightforward to check that the same bound holds also when one of the edge's endpoints is in $V_R$ or $V\setminus V_{CR}$. Thus, using the Ring Lemma's (Lemma 4.2) constant $A=A(\Delta)$ from part (i),
\begin{align*}
\mathcal E(h)
&\le\sum_{x\in V_{CR}\setminus V_R}\sum_{y:y\sim x}
\frac{((A+1)\operatorname{rad}(x))^2}{(C-1)^2R^2}\\
&\le\frac{\Delta(A+1)^2}{\pi(C-1)^2R^2}\sum_{x\in V_{CR}\setminus V_R}\operatorname{area}(C_x),
\end{align*}
where $\operatorname{area}(C_x)$ is the area that $C_x$ encloses (that is, $\pi\operatorname{rad}(x)^2$). We have that $\sum_x\operatorname{area}(C_x)\le\operatorname{area}(B(\mathbf0,2CR))=4\pi C^2R^2$, hence if $C=A+2$, then $\mathcal E(h)\le4\Delta C^2$, and the result follows for $c=(4\Delta C^2)^{-1}$.
\end{proof}
\begin{proof}[Proof of Theorem 4.4 (3)]
Consider the unit current flow $I$ from $v_0$ to $\infty$ and fix any $R\ge1$. Restricting this flow to the edges which have at least one endpoint in the annulus $V_{CR}\setminus V_R$ gives a unit flow from $V_R$ to $V\setminus V_{CR}$, by part (i) of Lemma 4.9. Hence, by part (ii) of that lemma and by Thomson's principle (Theorem 2.28), the energy contributed to $\mathcal E(I)$ from these edges is at least $c$.
In the same manner, the edges which have at least one endpoint in the annulus $V_{C^{2k+1}R}\setminus V_{C^{2k}R}$ contribute at least $c$ to $\mathcal E(I)$. Part (i) of Lemma 4.9 implies that all these edge sets are disjoint, hence $\mathcal E(I)=\infty$ and we learn that $G$ is recurrent (Corollary 2.39).
\end{proof}

\paragraph{Section 4.4.2: Proof of part 4.}
We will use the given circle packing of $G$ to create a random path to infinity with finite energy. This gives transience by Claim 2.46. This proof strategy is similar to that of Theorem 2.47.
\begin{proof}[Proof of Theorem 4.4 (4)]
Let $v_0$ be a fixed vertex of the graph, and apply a M\"obius transformation to make the circle of $P$ corresponding to $v_0$ be centered at the origin $0$.
We now use Claim 2.46 to construct a flow $\theta$ from $v_0$ to $\infty$ by choosing a uniform random point $p$ on $\partial\mathbb U$, taking the straight line from $0$ to $p$ and considering the set of all circles in the packing $P$ that intersect this line in the order that they are visited; this set forms an infinite simple path in the graph which starts at $v_0$.
To bound the energy of the flow, we claim that there exists some constant $C$ (which may depend on the graph $G$ and the packing $P$) such that the probability that the random path uses the vertex $v$ is bounded above by $C\operatorname{rad}(v)$. Indeed, since there are only finitely many vertices with centers at distance at most $1/2$ from $0$, we may assume that the center of $v$ is of distance at least $1/2$ from $0$. In this case, in order for $v$ to be included in the random path the circle of $v$ must intersect the line between $0$ and $p$. By the Ring Lemma (Lemma 4.2) the neighbors of $v$ have circles of radii comparable to $\operatorname{rad}(v)$ and so the probability of the line touching them is at most $C\operatorname{rad}(v)$. Since the vertex degree is bounded by $\Delta$ and $\sum_{v\in V}\pi\operatorname{rad}(v)^2$ is at most the area of $\mathbb U$, we find that
\[
\mathcal E(\theta)\le C\Delta\sum_{v\in V}\operatorname{rad}(v)^2\le C\Delta.
\]
Hence $G$ is transient by Corollary 2.39.
\end{proof}
\subsection*{整理说明（非作者正文）}
本题从已经给定的圆堆积出发，不要求另外建立所有无限图的圆堆积存在性；未把未收录的第(1)、(2)项算作题面结论。两个分项和环带电阻估计的证明均保留。电网络 Dirichlet/Thomson 原理、随机路径诱导流以及有限能量与暂留性的对应列为标准先修；Ring Lemma 保留作者原有几何论证。难点是选择截断径向函数获得各尺度一致的电阻下界，以及把圆盘中的随机射线化成能量可控的流。

圆心记号按原 PDF 的 $\operatorname{cent}$ 转录；网页公式中的索引排版以同版 PDF 可读文本为准。原图4.1只辅助 Ring Lemma 的相邻小圆讨论，未复制示意图；原文其他说明及式子均可按编号定位。取得的文本保存为转录稿，不称为 PDF 下载件。
\subsection*{可修改的简短标注（非作者正文）}
$c$：平面三角剖分上随机游走的常返性与暂留性。

$a$：圆堆积；环引理；Dirichlet 能量；有效电阻；Thomson 原理；随机路径诱导流；面积控制。

\texttt{cross\_theory\_status = yes}。圆的半径与面积分别控制环带上的电位差和随机路径的经过概率，再通过电网络能量判定返回随机游走类型。位置：Lemma 4.9、\S4.4.1--4.4.2。
标签为非穷尽、可修改初稿。
\subsection*{许可与修改记录}
原作者及作品见来源栏；来源采用 CC BY 4.0。本题证摘录和附加整理沿用该许可。2026-09-28：增加题号、来源定位、独立排版、整理说明和初稿标注；数学内容的取材范围及排版变化见整理说明。
\end{document}
